\documentclass[11pt]{article}

\usepackage[margin=1in]{geometry}
\usepackage[T1]{fontenc}
\usepackage[utf8]{inputenc}
\usepackage{microtype}
\usepackage{amsmath,amssymb}
\usepackage{booktabs}
\usepackage{graphicx}
\usepackage{xcolor}
\usepackage{hyperref}
\usepackage[nameinlink,capitalise]{cleveref}
\usepackage{enumitem}
\usepackage{natbib}
\setlength{\emergencystretch}{2em}

\hypersetup{
  colorlinks=true,
  linkcolor=blue!45!black,
  citecolor=green!35!black,
  urlcolor=blue!55!black
}

\newcommand{\systemname}{Preference Intelligence}
\newcommand{\evidencegap}[1]{%
  \begin{center}
  \fcolorbox{orange!70!black}{orange!4}{%
    \parbox{0.90\linewidth}{\textbf{Not measured in the current study:} #1}}
  \end{center}}

\title{User-Owned Preference Policies for Cross-Model, Domain-Conditioned
Language-Model Personalization}
\author{Anonymous authors}
\date{}

\begin{document}
\maketitle

\begin{abstract}
Language-model products commonly keep response preferences inside one provider
or conflate preferences with conversational memory. We formulate personalization
as a user-owned behavioral policy: typed preferences are selected by domain and
task, compiled into a provider-neutral instruction, and remain inspectable,
reversible, and subordinate to the current request. We contribute an auditable
representation with scope, confidence, provenance, and locks; five comparison
conditions; a controlled multidomain benchmark; and a reproducible execution and
blinded-evaluation pipeline. In an exploratory two-model local pilot comprising
128 real generations, a compact structured policy used 19.0 fewer input tokens
per prompt than a short history-retrieval baseline, but produced 70.8 more words
and added 3.12 seconds of latency on average. Its post-hoc reference-concept
coverage was indistinguishable within a wide exploratory interval
($\Delta=0.005$, 95\% cluster-bootstrap interval $[-0.141,0.177]$). The result
shows that input-context compression alone does not establish end-to-end
efficiency. No human ratings have yet been collected, so improved subjective
quality remains an open hypothesis rather than a reported result.
\end{abstract}

\section{Introduction}
\label{sec:introduction}

People use multiple large language models (LLMs), but each interface typically
receives only a local account setting, a manually repeated instruction, or a
provider-specific memory. The resulting problem is not only remembering what a
person said. It is representing \emph{how that person wants a response to be
constructed for the task at hand} and transferring that representation without
assuming that expertise or style in one domain applies everywhere.

We call this proposed layer \systemname{}. Its design goals are portability across
model providers, contextual selection, and transparent user control. A response
preference is an instruction-like behavioral constraint---such as desired depth,
format, or use of examples---rather than a demographic label or an assertion
about identity. The current request always outranks a stored default.

The primary research question is whether a dynamically learned,
domain-conditioned, provider-independent preference representation improves
subjective response quality across heterogeneous model families. We further ask
whether it improves on static profiles and history retrieval, whether learning is
reliable under drift and adversarial evidence, and whether structured profiles
reduce context overhead.

The contributions of the present artifact are:
\begin{enumerate}[leftmargin=*]
  \item a formal, inspectable \emph{preference policy} representation and
  applicability problem that separates response behavior from conversation memory;
  \item a behavioral-conformance framing of portability: an unchanged profile
  must induce the intended scoped behavior across model families, not merely
  serialize successfully;
  \item a benchmark and human-study protocol targeting domain isolation,
  preference learning, context efficiency, drift, and incorrect inference; and
  \item a local-first reference implementation and reproducible comparison of five
  personalization conditions, including 128 real local-model generations and an
  observed prompt-compression/output-expansion trade-off.
\end{enumerate}

The representation and pipeline are implemented contributions; the pilot is
exploratory systems evidence. It does not establish the primary subjective-quality
hypothesis. Irrelevant, stale, or incorrect preferences can still reduce quality
or undermine user control.

\section{Related Work}
\label{sec:related}

The relevant literature spans personalized language generation, user modeling,
implicit feedback, adaptive interfaces, recommender systems, conversational
memory, persona conditioning, and preference optimization. Context-aware
recommender systems and implicit-feedback methods established that behavior is
not equivalent to positive preference and that context belongs in the utility
model \citep{hu-etal-2008-collaborative,adomavicius-2011-context-aware}.

Recent LLM personalization work has moved beyond fixed persona sentences. LaMP
established retrieval-based personalization tasks
\citep{salemi-etal-2024-lamp}; other systems use per-user adapters or embeddings
\citep{tan-etal-2024-democratizing,liu-etal-2025-llms}, domain-aware compressed
profiles \citep{su-etal-2025-personalized}, and latent preference attributes
\citep{li-etal-2025-prefpalette}. PLUS jointly learns an interpretable natural-
language summary and a personalized reward model \citep{nam-etal-2026-plus}, while
S2Pref distinguishes stable from situational preferences and includes preference
conflicts \citep{gao-etal-2026-beyond}. These results mean that dynamic profiles,
domain conditioning, and interpretable summaries are not individually novel.

Evaluation work also warns against assuming that personalization helps. PrefEval
tests explicit and implicit preference following in long contexts
\citep{zhao-etal-2025-prefeval}, and PrefDisco reports settings where uninformed
personalization underperforms a generic response \citep{li-etal-2026-prefdisco}.
Broad memory systems offer temporal retrieval and compression mechanisms
\citep{park-etal-2023-generative,packer-etal-2023-memgpt,chhikara-etal-2025-mem0}
but usually optimize episodic or factual recall rather than an editable response-
preference policy.

The narrower gap investigated here is the joint system of typed response
dimensions, hierarchical applicability, calibrated evidence provenance, user
governance, provider-neutral compilation, and behavioral cross-model evaluation.
The structured literature review found no peer-reviewed demonstration of the
entire bundle at its evidence cutoff; this is an absence-of-evidence statement,
not a priority claim. The review protocol and study selection limitations are
reported in the accompanying artifact.

\section{Problem Formulation}
\label{sec:formulation}

Let $u$ denote a user, $x$ a current request, $d$ its domain, $t$ its task type,
$m$ a model provider/snapshot, and $P_u$ the portable preference profile. A
selector $S(P_u,d,t,x)$ returns a bounded set of relevant preference records,
while compiler $C$ maps those records to an instruction context $c$. Generation
is
\begin{equation}
  y \sim p_m\!\left(y \mid x, C(S(P_u,d,t,x))\right).
\end{equation}
The current request is a hard precedence constraint: if $x$ conflicts with a
stored default, $x$ wins. Selection is therefore not merely nearest-neighbor
retrieval; it combines scope compatibility, evidence authority, confidence,
recency, locks, and conflict resolution.

For subjective utility $Q(u,x,y)$, the primary estimand is the average
within-user contrast between domain-dynamic personalization and no
personalization over a prespecified population of users, tasks, and model
families. Portability additionally requires that $P_u$ is unchanged across $m$;
thin message-format adapters do not count as model-specific preference storage.
We therefore distinguish \emph{syntactic portability} (the profile can be parsed
and compiled for a new model) from \emph{behavioral conformance} (the resulting
response satisfies the same applicable policy without increasing correctness,
safety, privacy, or cost failures beyond a prespecified margin).

The design treats correctness, safety, privacy exposure, context cost, and
preference adherence as distinct outcomes. A system that increases stylistic
adherence while decreasing correctness does not unconditionally improve quality.

\section{Preference Representation}
\label{sec:representation}

Each preference record contains a stable identifier, bounded dimension and value,
scope tuple $(d,s,t)$ for domain, subdomain, and task, confidence, lifecycle state,
source type, evidence counts, timestamps, decay policy, user lock, user-authored
evidence, and provenance. Null scope components form a global default; more
specific compatible records override broader records for the same dimension.

This hybrid symbolic representation is human-readable and deterministically
filterable while allowing learned components to propose updates. Natural-language
compilation occurs only after selection. Assistant responses and webpage text are
untrusted inputs and cannot independently become user evidence. Users can inspect,
edit, lock, suppress, or delete values.

\evidencegap{Confidence calibration and observed editability require approved
participant data. Schema conformance is covered by software tests, but is not a
user outcome.}

\section{Dynamic Preference Learning}
\label{sec:learning}

The proposed updater aggregates signed, timestamped evidence by dimension and
scope. Explicit dashboard changes and direct statements have greater authority
than repeated implicit corrections; one ambiguous event cannot create a durable
high-confidence preference. A newer direct durable statement can supersede an
unlocked inferred value, while weaker contradictory evidence maintains an
explicit ambiguous state and causes selection to abstain. Locked values do not
change automatically, but the conflict remains inspectable.

Recency decay is applied to inferential evidence, while user-locked preferences do
not decay. Temporary, task-local requests affect the current generation but are
not promoted to durable profile values without an explicit persistence cue. Drift
is evaluated through stable change points, temporary overrides, and alternating
task-dependent preferences.

The reference interface also accepts negative applicability evidence: ``this
preference was not relevant here'' adds a scoped exception while retaining the
underlying preference assertion. This creates direct supervision for the second
error type without conflating it with preference truth.

\evidencegap{Learning curves, activation precision/recall, calibration, and
adaptation-delay estimates require the frozen longitudinal experiments.}

\section{Domain-Conditioned Personalization}
\label{sec:conditioning}

Domain conditioning prevents a local observation---for example, advanced Java
knowledge---from becoming a general assumption about finance or physics. The
selector admits global defaults and compatible domain, subdomain, and task
records, resolves each dimension by specificity and authority, and applies a
confidence threshold. Missing domain knowledge remains unknown.

The reference implementation operationalizes a separate applicability gate as a
weighted combination of symbolic scope match, dimension--task relevance, and
decayed evidence confidence. Scope mismatch is a hard rejection; locked state
does not bypass it. The present coefficients are inspectable engineering priors,
not fitted probabilities or experimental findings. Every application and
rejection produces a bounded explanation trace for error analysis.

The controlled benchmark deliberately pairs high familiarity in one domain with
low familiarity in another. It also includes ambiguous domains and explicit
current-request overrides. The central mechanism ablation removes scope while
holding the evidence and token budget fixed.

\evidencegap{Domain leakage, relevant-preference recall, and domain-conditioned
versus global contrasts require the locked generation and blinded rating study.}

\section{Experimental Setup}
\label{sec:setup}

We predefine five conditions: no personalization (A), a participant-authored
static global profile (B), retrieval of relevant user-authored history (C), a
dynamic global structured profile (D), and a dynamic domain-conditioned profile
(E). Generation settings and neutral task instructions are held fixed within a
comparison. Equal personalization-token budgets support the primary quality
comparison; natural budgets support the efficiency comparison.

The synthetic benchmark contains multidomain personas, development and locked
task-template families, explicit conflicts, invalid evidence, and preference
drift. The proposed human evaluation is a randomized, blinded,
repeated-measures incomplete block with participant and task effects. The primary
outcome is pairwise E-versus-A preference with ties retained. Secondary outcomes
include usefulness, clarity, depth/detail appropriateness, adherence, context
tokens, latency, domain leakage, and incorrect activation.

\paragraph{Reproducibility.} Every generation record includes the condition,
provider and returned model identifiers, parameters, full prompts, selected
preference records, compiled context, token usage or labeled approximation,
latency, response, automatic diagnostics, error, timestamp, and seed. Run
manifests contain input hashes and the code revision when available.

\paragraph{Executed local-model pilot.} We instantiated one synthetic persona with
advanced software-engineering and beginner-finance preferences and selected four
tasks: Java virtual threads, Kafka rebalancing, bond duration, and portfolio
diversification. Qwen~3 8B Q4\_K\_M and Llama~3.2 1B Q8\_0 were served through a
local Ollama endpoint. We disabled hidden reasoning to reserve a comparable
260-token visible-output budget, set temperature to 0.3, generated two stochastic
replicates per cell, and deterministically shuffled cells using recorded seeds.
The first run crossed both models with all five conditions ($80$ generations;
seed 20260913). After observing excess structured-prompt overhead, we retained
that run and executed a labeled post-hoc compiler ablation over A, C, and E
($48$ generations; seed 20260914). No failed cells were discarded.

\paragraph{Pilot analysis.} A narrow reference-concept coverage diagnostic was
derived after generation from factual reference points already present in the
benchmark. It can miss valid paraphrases and is not a factuality or preference
metric. We report paired differences and exploratory percentile intervals from
5,000 cluster-bootstrap draws over eight model--task clusters. Because the metric
and compact ablation were post-hoc, we report no confirmatory $p$-values.

\paragraph{Human-study status.} No participants have been recruited and no human
ratings are available. Exact hosted-provider snapshots, participant count, and
smallest effect of interest must be frozen in a preregistration and approved study
protocol before confirmatory collection.

\section{Exploratory Local-Model Results}
\label{sec:results}

All 128 assigned local-model generations completed without a transport failure.
The primary five-condition run contributed 80 outputs and the compact-compiler
ablation contributed 48. Raw responses, complete prompts, selected preference
records, provider token counts, finish reasons, latency, run manifests, and seeds
are retained in JSONL. The two analysis datasets also provide 64 and 32
condition-blinded response pairs, respectively, with mapping keys stored separately.

\subsection{Five-condition pilot}

The initial structured compiler (\texttt{verbose\_v1}) did not deliver the hoped-for
context advantage. Compared with history retrieval, E consumed 65.0 additional
input tokens per prompt (exploratory 95\% interval $[62.1,68.0]$), produced 52.1
additional words ($[11.6,99.1]$), and added 3.56 seconds of latency
($[1.16,6.26]$). Its reference-concept coverage difference was $-0.047$
($[-0.177,0.104]$). Compared with no personalization, E produced 34.8 fewer words
($[-75.2,-7.0]$) and returned 1.13 seconds faster ($[-2.29,-0.13]$), but required
116.0 more input tokens. The coverage difference was $0.073$
($[-0.052,0.188]$). These intervals describe this small post-hoc pilot and are not
confirmatory confidence intervals for a target population.

\subsection{Compact compiler ablation}

We replaced descriptive structured records with fixed, bounded behavioral clauses
(\texttt{compact\_v2}), reducing the example compiled profile from 447 to 158
characters. The resulting E prompts used 19.0 fewer input tokens than history
retrieval ($[-19.75,-18.25]$), thereby meeting a narrow input-context objective.
However, E produced 70.8 more words ($[28.6,117.3]$), 99.1 more output tokens
($[42.0,161.8]$), and added 3.12 seconds of latency ($[1.32,5.03]$). The rate of
length-limited termination was 0.313 higher, although its interval included zero
($[-0.063,0.750]$). Reference-concept coverage was nearly unchanged
($\Delta=0.005$, $[-0.141,0.177]$).

Against no personalization, compact E used 32.0 additional input tokens
($[31.25,32.75]$), produced 16.7 fewer words ($[-42.8,6.3]$), and returned 1.47
seconds faster ($[-2.73,-0.47]$). Its coverage difference was $-0.021$
($[-0.156,0.125]$).

\begin{table}[t]
  \centering
  \small
  \caption{Pooled paired E-minus-comparator differences across two models, four
  tasks, and two replicates. Intervals are exploratory model--task cluster
  bootstraps. Lower is better for token, word, and latency costs; reference
  coverage is a post-hoc lexical proxy.}
  \label{tab:pilot-results}
  \begin{tabular}{llrrr}
    \toprule
    Compiler & Comparator & $\Delta$ coverage & $\Delta$ input tokens &
    $\Delta$ words \\
    \midrule
    verbose\_v1 & History retrieval & $-0.047$ & $+65.0$ & $+52.1$ \\
    verbose\_v1 & No personalization & $+0.073$ & $+116.0$ & $-34.8$ \\
    compact\_v2 & History retrieval & $+0.005$ & $-19.0$ & $+70.8$ \\
    compact\_v2 & No personalization & $-0.021$ & $+32.0$ & $-16.7$ \\
    \bottomrule
  \end{tabular}
\end{table}

The strongest supported statement is therefore a systems trade-off: compact
structured selection reduced input context relative to this history baseline, but
the tested instruction semantics expanded generation enough to erase the end-to-end
efficiency benefit. Neither run establishes preference satisfaction or subjective
quality.

\section{Ablation Studies}
\label{sec:ablations}

The broader planned ablations independently remove domain scope, confidence gating,
recency decay, evidence-source weighting, provenance filtering, conflict handling,
and user locks. Representation baselines compare structured records with a
natural-language profile and multiple history retrievers. Tests hold evidence,
task allocation, generation settings, and applicable token budget fixed.

The first completed ablation changed only the compiler surface form from a verbose
record serialization to compact behavioral clauses. It reduced the representative
compiled profile from 447 to 158 characters and reversed the input-token contrast
against history retrieval from $+65.0$ to $-19.0$ tokens per prompt. It did not
improve the exploratory reference proxy and increased both generated length and
latency (\cref{sec:results}). This identifies compiler wording and total generated
tokens---not profile storage alone---as necessary experimental factors.

\evidencegap{The remaining preregistered component ablations and their
full-versus-ablation effects with corrected intervals,
including adherence precision and recall, false activations, domain leakage,
adaptation delay, context cost, and safety failures. An ablation that activates
every preference must not be credited for recall without its precision cost.}

\section{Human Evaluation}
\label{sec:human}

The human-study protocol separates target-user preference judgments from external
annotation of correctness and constraints. Participants rate debranded responses
before seeing preference explanations. A later transparency task measures whether
they can understand, edit, lock, delete, and disable inferred values. Recruitment
cannot begin until the applicable institutional review determination is complete.

\evidencegap{An approved protocol identifier, recruitment and compensation
details, sample flow, preregistered exclusions, rating distributions, reliability
before adjudication, blinding assessment, and qualitative coding. No human data
have been collected.}

\section{Discussion}
\label{sec:discussion}

The pilot supports syntactic portability but not yet behavioral conformance: one
unchanged structured profile was accepted by two heterogeneous model families,
yet the small automatic analysis cannot show that either family produced answers
the target user would prefer. That distinction prevents a successful serializer
or prompt injector from being mislabeled as successful personalization.

The compiler ablation also exposes a cost-accounting trap. A compressed policy can
save input tokens while increasing generated tokens and wall-clock latency. In our
sample, directive-like compact clauses appear to have encouraged fuller responses
despite a global concise preference. This interpretation is mechanistic and
post-hoc; output caps and unequal model sizes are alternative explanations. Future
work should jointly optimize adherence, correctness, input tokens, output tokens,
and truncation rather than treating prompt length as the efficiency endpoint.

The near-zero compact E-versus-history difference on reference coverage should not
be read as equivalence. The interval is wide, the diagnostic is lexical, and only
one synthetic persona and four tasks were used. Conversely, the absence of a
positive automatic result is useful counterevidence against assuming that typed
profiles automatically improve answers. The condition-blinded files now make the
next decisive step concrete: target-user preference judgments and independent
correctness review under a frozen analysis plan.

The larger research program should distinguish broad benefit, style-only benefit,
provider and domain heterogeneity, and net harm from incorrect or stale preferences.
Practical effect sizes, correction burden, and abstention quality matter more than
statistical significance alone.

Implementation of contextual negative feedback further sharpens the research
program. Asking whether a stored preference is true and asking whether it belongs
in a current prompt are separable labeling tasks. Future human evaluation should
therefore report inference calibration, applicability precision, abstention
coverage, and recovery after each error type rather than collapsing them into one
personalization score.

\section{Privacy, Safety, and Limitations}
\label{sec:limitations}

A portable profile can reduce repeated disclosure but can also centralize
sensitive behavioral information. The proposed system is local-first, minimizes
raw conversation retention, records provenance, rejects assistant/webpage text as
preference authority, and makes profile values user-visible and deletable.
Cross-provider use remains a disclosure event and must be explicit and scoped.

Preferences never authorize unsafe output. The representation avoids unnecessary
sensitive traits and stores desired explanation behavior rather than demographic
inferences. Synthetic privacy canaries are used for leakage tests; real secrets
must not be planted.

The benchmark's controlled personas may exaggerate coherence and stability,
automatic checks reward surface compliance, English-only tasks limit cultural and
linguistic generalization, provider updates challenge reproducibility, and
self-reported familiarity is not a complete expertise measure. Repeated exposure
can construct preferences rather than merely reveal them. The history-retrieval
baseline is sensitive to retriever quality. A positive average effect could hide
provider-specific or domain-specific harm.

The executed pilot adds several concrete limitations. It uses two quantized local
models of very different capacity, one synthetic persona, four English tasks, two
replicates, a 260-token cap, and a post-hoc lexical concept diagnostic. Several
outputs stopped at the length limit. The local endpoint avoids third-party data
disclosure but does not represent hosted-provider policies or behavior. Cluster
bootstrap intervals over only eight model--task clusters are unstable. The
compiler ablation was motivated after inspection of the first run and is therefore
hypothesis-generating. No human preference, editability, privacy-leakage, or
longitudinal-learning outcome was measured.

\evidencegap{Adverse-output, privacy-exposure, subgroup, accessibility, and hosted
provider data-retention outcomes. The local pilot contained no participant data.}

\section{Conclusion}
\label{sec:conclusion}

This work frames response personalization as a portable, contextual, transparent,
and user-governed policy-selection problem rather than undifferentiated long-term
memory. It contributes a typed representation, behavioral-conformance definition,
benchmark, executable five-condition harness, local-first product, and blinded
human-evaluation path.

Across 128 real generations from two local model families, the same profile and
selection pipeline executed without failures. A compact compiler saved 19 input
tokens per prompt relative to a short history baseline, but produced about 99 more
output tokens and added 3.12 seconds of latency; its exploratory reference coverage
was not distinguishable in this small sample. Thus the measured contribution is
an auditable systems artifact and a falsifiable efficiency trade-off, not evidence
that personalization improves subjective quality. Establishing that claim requires
the preregistered target-user study and independent correctness evaluation.


\section*{Reproducibility statement}
The repository contains versioned benchmark inputs, deterministic mock generation,
provider adapters, condition compilation, randomized real local-model runs, raw
JSONL responses, immutable input hashes, analysis code, and condition-blinded
annotation files. The reported pilot data use Qwen~3 8B and Llama~3.2 1B served
locally; all 128 assigned generations completed. No human judgments, private
participant data, API credentials, or fabricated outcomes are included.

\section*{Acknowledgments}
Withheld for anonymous review.

\IfFileExists{../literature/references.bib}{%
  \bibliographystyle{plainnat}
  \bibliography{../literature/references}
}{%
  \paragraph{Bibliography status.} A verified bibliography has not yet been
  linked. Do not insert placeholder citations.
}

\end{document}
